%! AP Statistics — Unit 5, Lesson 5.1 — Homework (Answer Key)
\documentclass[10pt]{article}
\usepackage{apstats-article}
\usepackage{apstats-key}   % pulls in -boxes; do NOT also load -boxes
\usepackage{grading}
\gradingassignment{apstats-u05-l01-hw}
\gradingcourse{advanced_statistics}
\newcommand{\TallMath}[1]{$\displaystyle #1\rule[-1.4em]{0pt}{3.2em}$}

\begin{document}

\pageheader{Unit 5, Lesson 5.1}{Homework --- Answer Key}
{\small\itshape Write each final answer in its box. Only what is inside a box is graded.}\par

\vspace{0.15in}

\begin{enumerate}

  \item
        \begin{enumerate}[label=\alph*.]
          \item
                \begin{itemize}
                  \boxguard[4]
                  \item Statistic: \ans{sample mean heart rate, 72 bpm}\gradedkey{\gradedblank[id=hw1a1, type=short, points=1, answer={sample mean 72}, check=text, display={\bfseries sample mean 72}]{3in}{0.3in}}
                        \par Notation: \ans{$\bar{x}$}\gradedkey{\gradedblank[id=hw1a2, type=short, points=1, answer=\bar{x}, check=expression, display={\boldmath$\bar{x}$}]{0.5in}{0.19in}}
                  \boxguard[4]
                  \item Parameter: \ans{true mean heart rate of all adults}\gradedkey{\gradedblank[id=hw1a3, type=short, points=1, answer={mean heart rate of all adults}, check=text, display={\bfseries mean of all}]{3in}{0.3in}}
                        \par Notation: \ans{$\mu$}\gradedkey{\gradedblank[id=hw1a4, type=short, points=1, answer=\mu, check=expression, display={\boldmath$\mu$}]{0.5in}{0.19in}}
                \end{itemize}
          \item
                \begin{itemize}
                  \boxguard[4]
                  \item Statistic: \ans{sample proportion who exercise $\ge$ 3 times/week, $64/200 = 0.32$}\gradedkey{\gradedblank[id=hw1b1, type=short, points=1, answer={sample proportion 0.32}, check=text, display={\bfseries sample prop. 0.32}]{3in}{0.3in}}
                        \par Notation: \ans{$\hat{p}$}\gradedkey{\gradedblank[id=hw1b2, type=short, points=1, answer=\hat{p}, check=expression, display={\boldmath$\hat{p}$}]{0.5in}{0.19in}}
                  \boxguard[4]
                  \item Parameter: \ans{true proportion of all college students who exercise $\ge$ 3 times/week}\gradedkey{\gradedblank[id=hw1b3, type=short, points=1, answer={proportion of all college students who exercise 3+ times a week}, check=text, display={\bfseries prop. of all}]{3in}{0.3in}}
                        \par Notation: \ans{$p$}\gradedkey{\gradedblank[id=hw1b4, type=short, points=1, answer=p, check=expression, display={\boldmath$p$}]{0.5in}{0.19in}}
                \end{itemize}
        \end{enumerate}

  \boxguard[22]
  \item
        \begin{enumerate}[label=\alph*.]
          \item \ansline{No --- they are two different sample means from two different randomly selected groups of seedlings, so they are not the same measurement. Each describes only the 15 seedlings in that sample.}\gradedkeyline{\gradedblank[id=hw2a, type=rubric, points=2, display={No: two different samples of 15 seedlings; each mean describes only its own sample.}]{\linewidth+6pt}{\gradedlinesheight{1}}}
          \boxguard[17]
          \item \ansline{No. Different random samples from the same greenhouse will naturally produce different sample means by chance (sampling variability). Neither measurement is wrong; they reflect natural variation.}\gradedkeyline{\gradedblank[id=hw2b, type=rubric, points=2, display={No: different random samples give different means by chance (sampling variability).}]{\linewidth+6pt}{\gradedlinesheight{1}}}
          \item \ans{sampling variability}\gradedkey{\gradedblank[id=hw2c, type=short, points=1, answer={sampling variability}, check=text]{1.5in}{0.17in}}
        \end{enumerate}

  \boxguard[8]
  \item
        \begin{enumerate}[label=\alph*.]
          \item \ans{the sampling distribution of the sample mean}\gradedkey{\gradedblank[id=hw3a, type=short, points=1, answer={sampling distribution of the sample mean}, check=text]{2.1in}{0.17in}}
          \boxguard[17]
          \item \ansline{Yes. Because the samples are random and come from a population with mean 100, the sample means should vary around 100 on average. The center of the sampling distribution equals the population mean.}\gradedkeyline{\gradedblank[id=hw3b, type=rubric, points=2, display={Yes: sample means vary around $\mu = 100$; the sampling distribution's center equals the population mean.}]{\linewidth+6pt}{\gradedlinesheight{1}}}
          \boxguard[17]
          \item \ansline{The spread (standard deviation) of the dot-plot would decrease. Larger samples reduce sampling variability because averaging more values smooths out the extremes.}\gradedkeyline{\gradedblank[id=hw3c, type=rubric, points=2, display={The spread would decrease: larger samples mean less sampling variability.}]{\linewidth+6pt}{\gradedlinesheight{1}}}
        \end{enumerate}

\end{enumerate}

\begin{extensionbox}
\textbf{Extension (optional).} \textit{Responses will vary. Look for: (1) identification of the two statistics and the claimed parameter, (2) acknowledgment that sampling variability alone can produce differences without a real population difference, and (3) some reasoning about sample size or context.}
\end{extensionbox}

\newpage % the blank's reflection box falls on page 3; keep the key page for page
\begin{reflectionbox}
\textbf{Preview:} Lesson 5.2 revisits the normal distribution, which will be the key
model for describing the shape of sampling distributions. Review your notes on
z-scores and normal probabilities from Unit 1.
\end{reflectionbox}

\end{document}
